\chapter{The console}\label{ch:console}
\begin{goals}
\textbf{In this chapter} you drive wuffle by typing: its commands, the shell around them, and routines
you can save and run again.
\end{goals}

Everything you do with the mouse on the desk, you can also type. The console is a real shell,
\code{geas}, and wuffle's commands are part of it, so they work with pipes, loops and files like any
other command. A sequence of commands saved in a file is a \term{routine}: run it again tomorrow, on
another project, and it does the same thing.

\shot{console}{The console after \code{sample tocher}, \code{contour kamb 3} and \code{stats}: the commands
draw on the net and answer in the shell.}

\section{Opening it}

Press \key{`} (backquote), choose \menu{View>Console}, or click the \ui{console} tab under the diagrams.
The prompt shows the working directory, which starts at \file{/project}: the project itself, as files
(\cref{ch:project-fs}).

\key{↑} and \key{↓} recall earlier lines. \key{Tab} completes: at the start of a line it completes a
command; after one, a dataset, layer or diagram name, then a path. A name with spaces comes back quoted.

\section{wuffle's commands}\label{sec:commands}

Commands act on the \emph{active diagram}, the tab in front. A layer is named by its own name or its
dataset's, and the start of the name is enough (\code{stats toch}). With no name, a command takes the
first visible layer.

\begin{center}\small
\begin{tabularx}{\linewidth}{@{}>{\ttfamily\raggedright}p{0.36\linewidth}L@{}}\toprule
\normalfont\sffamily\bfseries command & \sffamily\bfseries what it does \\\midrule
sample tocher & open the Tocher sample \\
diagrams & list the diagrams; \code{●} marks the active one \\
use \emph{name} & make a diagram active \\
new net|rose|fabric & a new diagram showing the active one's layers \\
layers & the active diagram's layers, \code{☑} visible \\
show|hide \emph{element} [\emph{layer}] & poles, circles, contours, mean, eigen, labels, v1, v2, v3 \\
contour kamb [\emph{k}] | fisher & contours on, by Kamb's method \parencite{kamb1959} at \emph{k}\,$\sigma$, or with a Fisher kernel \\
select cone \emph{att} \emph{angle} & the poles within \emph{angle} of a direction (a plane's pole, or a line) \\
select band \emph{plane} \emph{width} & the poles within \emph{width} of a great circle \\
select all|none|invert & on every layer of the active diagram \\
stats [\emph{layer}] & n, the Fisher mean with $\kappa$ and $\alpha_{95}$ \parencite{fisher1953}, the eigenvalues, and Woodcock's K and C \parencite{woodcock1977} \\
mean|eigen|beta [\emph{layer}] & keep the result as a dataset, drawn on this diagram \\
intersect \emph{a} \emph{b} & where two layers' mean planes meet, as a dataset \\
rotate \emph{tr→pl} | off & centre the net on a direction, or back to normal \\
undo | redo & the last change, whoever made it (\cref{sec:undo}) \\
verify [\emph{name}] & check field originals' signatures (\cref{sec:field-logs}) \\
theme dark|light|system & the colour theme \\
history & the lines typed so far \\
help [\emph{command}] & all of this, or one command \\
\bottomrule\end{tabularx}
\end{center}

A selection command adds to the selection with \code{+} and takes from it with \code{-}:
\code{select + cone 150/60 20}. Attitudes are typed as everywhere else: \code{150/60} a plane,
\code{150→60} a line.

\begin{caution}
In the console, \code{>} is the shell's: it sends output to a file. \code{select cone 150>60 10} is
read as \enquote{\code{select cone 150}, output into a file called \file{60}}, and in \file{/home} it
would make that file. Type the arrow, or quote the attitude: \code{select cone '150>60' 10}. (The
compass's typed entry, which is not a shell, takes \code{150>60} as it is.)
\end{caution}

\begin{tryit}
\begin{enumerate}
\item Type \code{sample tocher}, then \code{contour kamb 3}. The net is contoured at 3\,$\sigma$.
\item Type \code{select cone 040/50 20}. The poles within 20° of the pole to \code{040/50} turn violet;
the console says how many.
\item Type \code{mean}. A new dataset, the mean of the layer, is drawn on the net, and the tree lists it
with an arrow: it remembers it was derived.
\end{enumerate}
\end{tryit}

\section{The shell around them}\label{sec:shell}

Besides wuffle's commands the console has a Unix toolbox: \code{ls}, \code{cat}, \code{head},
\code{tail}, \code{grep}, \code{sort}, \code{uniq}, \code{wc}, \code{cut}, \code{tr}, \code{find},
\code{xargs}, \code{tee}, \code{mkdir}, \code{mv}, \code{cp}, \code{rm}, \code{ln}, \code{tar},
\code{zip} and more; \code{help} lists them all. There is no \code{sed} or \code{awk} yet. Pipes,
\code{\$(…)}, variables, \code{for} and \code{while} loops, \code{if} and functions work the way they do
in a Unix shell.

wuffle's commands write to the pipe like the rest:

\begin{lstlisting}
stats tocher | grep mean
for m in fisher kamb; do contour $m; done
wc -l data/tocher.csv
\end{lstlisting}

Two places are yours. \file{/home} is saved with the project; \file{/tmp} is scratch, gone when the page
closes. Anything else outside \file{/project} is the shell's own.

\section{Routines}\label{sec:routines}

A routine is a shell script. Write it into \file{/home} and run it with \code{.} (dot):

\begin{lstlisting}
printf 'sample tocher\ncontour kamb 3\nshow mean\nstats\n' > /home/fabric.sh
cat /home/fabric.sh
. /home/fabric.sh
\end{lstlisting}

\code{history} prints what you have typed this session, which is the quickest way to start one. Because
\file{/home} travels in the project file, the routine goes wherever the project goes.

\begin{note}
A failed command stops nothing by itself: as in any shell, the script carries on. Chain steps with
\code{\&\&} when a later one makes no sense without an earlier one.
\end{note}
